% s7: the room grows (point 5, first half). ASCII only.
\section{The room grows}

\spine{5} The infinite world is the last room. The radius $R$ belongs to the parameter space of the family. Over each $R$ sits the
whole quantum theory in that room, with its carriers and its gap $\Delta(R)$, the lowest energy above the vacuum over the
entire vacuum-orthogonal physical Hilbert space (every charge, every parity). Changing the renormalization scale
to $1/R$ follows the running coupling, but changing $R$ also changes the physical geometry. These are related, not
identical operations. Asymptotic freedom controls small $R$, not a globally monotone strong-coupling law.

\textbf{Dimensional transmutation.} Pure Yang--Mills has no mass in its Lagrangian. The running manufactures one: holding
the theory fixed while the cutoff is removed defines a scale $\Lambda$ (from the coefficients $b_0=22/3$, $b_1=136/3$).
At fixed squashing $t$, group and $\theta=0$, a finite-room mass has the form $\Lambda f(R\Lambda;t)$.
A constant mass ratio belongs to the infinite-volume limit, if it exists; its value uses a fixed scheme for $\Lambda$.

\begin{claim}{PROVED}{The scaling reduction}\label{b:scaling}
Given that each continuum room exists, $H_R\cong R^{-1}\mathcal K(x)$ for a one-parameter family $\mathcal K(x)$ on the
unit sphere, $x=R\Lambda$. With $\delta(x)=R\,\Delta(R)$ the gap of $\mathcal K(x)$, $\Delta(R)/\Lambda=\delta(x)/x$, and
my uniform-gap route asks for $\inf_{x>0}\delta(x)/x>0$. This is stronger than the infinite-volume step alone needs:
a positive liminf along a converging family as $x\to\infty$. Dimensional analysis does not construct the rooms.
\evidence{Exact dimensional analysis; the existence of each room is H4$'$ (\S8)~\cite{ClosedSet}.}
\lines{20, 23}
\end{claim}

So the gap is a \emph{section} over the half-line $x\in(0,\infty)$, a choice of one number over each room, and the
infrared problem is that it cannot approach zero at the infinite end. Pointwise positivity is not enough.
My uniform bound has three parts.
\emph{(1) Continuity along the base.} A change of carrier between different charges is an exact crossing
(Statement~\ref{b:cross}), and the minimum of finitely many isolated continuous low-level curves is continuous, so such a switch never breaks the section; it
only changes who carries it. \emph{(2) No zero in any room:} a discrete spectrum and a single vacuum give $\delta(x)>0$.
With (1), compactness gives a positive minimum on every closed stretch away from the endpoints.

\begin{claim}{CONJECTURE}{Continuity and positivity on every compact stretch}\label{b:compact}
For every $0<x_0<X$, $\min_{[x_0,X]}\delta(x)/x>0$. Sub-conjectures: (i) compact resolvent and a unique ground state
survive the ultraviolet limit in each room; (ii) $\mathcal K(x)$ is norm-resolvent continuous in $x$~\cite{Kato}.
\evidence{Both hold for every lattice regularization (positivity-improving transfer matrix~\cite{OsterwalderSeiler78})
and for the slow sector by Theorem A; the divergence to remove is the one matched in \S6.}
\lines{20, 17, 19}
\end{claim}

\emph{(3) Control at both ends.} As $x\to0$, $g^2(R)\sim8\pi^2/(b_0\log(1/x))\to0$, and, under H2, the low spectrum is the reduced model
dressed, and $\delta(x)=\delta_0+O(g^2)$ with $\delta_0>0$: tiny rooms are enormously gapped in units of $\Lambda$.

\begin{claim}{CONDITIONAL}{The small-room end}\label{b:smallend}
$\delta(x)/x\to\infty$ as $x\to0$, given a positive limiting reduced gap and H2's exclusion of additional low states.
H8 conjectures the charge-0 carrier, but only $\delta_0>0$ is used here; its ordering need not be settled first.
\evidence{Conditional on H2 (Statement~\ref{b:m3}) in a strengthened form (bounds on the high block and the off-diagonal block);
the small-volume picture of~\cite{Luscher83,KollerVanBaal88,vanBaal00} moved to the sphere.}
\lines{16, 18, 20}
\end{claim}

The large end is the wall. Strong coupling is not it (\S3): every room's continuum limit sends the bare coupling to zero.
What the large end needs is a \emph{finite-size criterion}, the statistical-mechanics idea that mixing checked at one
finite scale forces it at every scale~\cite{DobrushinShlosman85,MartinelliOlivieri94}.

\begin{claim}{CONJECTURE}{The large-room end: one-scale mixing (MS)}\label{b:ms}
There are a physical scale $\rho_*$, an integer $L$ and $\alpha<1$ such that, in a 4-cube of side $L\rho_*$ tiled by
blocks of side $\rho_*$, the total influence of the boundary blocks on smoothed gauge-invariant observables of the central
block is at most $\alpha$, uniformly in boundary condition and cutoff. Then correlations decay at rate
$\log(1/\alpha)/(L\rho_*)$ and $\liminf_{x\to\infty}\delta(x)/x>0$.
\evidence{Proved for lattice spin systems~\cite{DobrushinShlosman85,MartinelliOlivieri94} and lattice Yang--Mills at
strong coupling~\cite{ShenZhuZhu23}; no version for gauge fields uniform in the spacing is known to me.}
\lines{20}
\end{claim}

\begin{claim}{CONJECTURE}{H3 (M2): the gap in the field's own scale}\label{b:h3}
There is $c>0$, depending only on the group and on $t\ne1$, with $\Delta(R)/\Lambda\ge c$ for every $R$. The gap is on
the orthogonal complement of the chosen vacuum, so this hypothesis also excludes other zero-energy states.
Route: Statements~\ref{b:scaling}--\ref{b:ms}.
\evidence{Lattice spectroscopy shows a gap that is a fixed multiple of $\Lambda$~\cite{MorningstarPeardon99}; step
scaling tracks the running across volumes~\cite{LWW91}. This is the infrared half of the problem~\cite{JaffeWitten}.}
\lines{20}
\end{claim}

Every room can be gapped and the theory still massless: if $\delta(x)\sim x^p$ with $0<p<1$, then $R\Delta(R)\to\infty$
while $\Delta/\Lambda\to0$, which only the ratio to $\Lambda$ sees. In GIGI I can follow the section room by
room, a switch showing as an exact crossing: evidence, in the spirit of step scaling~\cite{LWW91}, not proof. The rooms, as they grow, must finally become all of
space.
